
Data attribution methods estimate training example influence, but their performance across different model architectures remains underexplored. This work presents a systematic comparison of TRAK, EK-FAC, and TracIn across encoder-only (BERT) and decoder-only (GPT-2) transformers at matched conditions (~110-125M parameters, 12 layers, SST-2 mislabeled detection with 5\% injected label noise). The experiments reveal architecture-dependent patterns: TRAK exhibits architecture-invariance (<1\% AUC difference at all compute budgets), while TracIn shows a directional BERT advantage (~3\% at low compute, though not statistically significant with 2 seeds). EK-FAC, contrary to theoretical predictions that Kronecker factorization better fits causal attention structure, shows no meaningful decoder advantage (0.27\% difference, p=0.90). The causal mechanism is validated through intermediate measurements: attention structure differs by 98.82\% sparsity between architectures (BERT upper-triangle sparsity: 1.18\%, GPT-2: 100\%), creating an $11\times$ Hessian eigenvalue difference (GPT-2 top eigenvalue: 0.502, BERT: 0.046). For practitioners, TRAK provides a robust default for cross-architecture attribution work; TracIn may offer advantages in encoder-only pipelines; EK-FAC performs equivalently on both architectures.
